% Generated draft: requires mathematical and source review.
\documentclass[11pt]{article}
\usepackage[margin=25mm]{geometry}
\usepackage{fontspec}
\setmainfont{Latin Modern Roman}
\usepackage{amsmath,amssymb,mathtools,unicode-math}
\setmathfont{Latin Modern Math}
\usepackage{enumitem,needspace,xurl,hyperref}
\setlength{\parindent}{0pt}
\setlength{\parskip}{5pt}
\setlength{\emergencystretch}{3em}
\newcommand{\N}{\mathbb{N}}
\newcommand{\Z}{\mathbb{Z}}
\newcommand{\Q}{\mathbb{Q}}
\newcommand{\R}{\mathbb{R}}
\newcommand{\C}{\mathbb{C}}
\newcommand{\F}{\mathbb{F}}
\newcommand{\A}{\mathbb{A}}
\newcommand{\PP}{\mathbb{P}}
\newcommand{\E}{\mathbb{E}}
\newcommand{\eps}{\varepsilon}
\newcommand{\dd}{\,\mathrm{d}}
\begin{document}
\section*{UnsolvedMath ID: 20001988}
\textbf{Source problem code:} catalogue entry
\subsection{Definitions and mathematical statement}
Let \(k\) be an algebraically closed nonarchimedean valued field and
\[
 R=k[x_1^{\pm1},\ldots,x_d^{\pm1}],\qquad
 \mathfrak p\subset R\text{ a proper prime ideal},\qquad
 r=\dim(R/\mathfrak p).
\]
Write \(\operatorname{Trop}(\mathfrak p)\subset\mathbb R^d\) for its nonarchimedean amoeba, equivalently tropicalization, using the coefficient valuation in initial ideals. Fix a finite pure \(r\)-dimensional polyhedral structure on it. Its simple facet graph \(G\) has one vertex per \(r\)-face, with an edge when two facets share an \((r-1)\)-face. Item A.6 asks whether \(G\) is always connected, a property called connectedness through codimension one. It does not ask for higher homotopy connectedness of the support.

The source background attributes ordinary connectedness to a then-forthcoming paper of Einsiedler, Lind and Kapranov. It also calls the amoeba homogeneous. For a general nontrivial coefficient valuation, a polyhedral complex need not be a homogeneous fan; that wording needs the appropriate constant-coefficient or valuation convention and is not imposed on the general statement here.

The aid reports an affirmative solution by Cartwright--Payne: the tropicalization of the irreducible variety defined by \(\mathfrak p\) is connected through codimension one, in every characteristic under the algebraically closed-field hypothesis. The zero-dimensional case is the single-point, single-facet case and is connected without any ridge adjacency. A single higher-dimensional facet is likewise connected as a one-vertex graph.

In characteristic zero the aid additionally invokes Maclagan--Yu. Let \(\ell\) be the lineality dimension in the chosen pure polyhedral structure, and put \(q=r-\ell\). The reported stronger graph consequence is that deletion of at most \(q-1\) facet vertices preserves connectedness, with the standard size qualifications. In particular, when
\[
 q>0,\qquad |V(G)|\geq q+1,
\]
the candidate corollary is
\[
 \kappa(G)\geq q=r-\ell,
\]
where \(\kappa\) is vertex connectivity. Menger's theorem then supplies \(q\) pairwise internally vertex-disjoint facet chains between any two distinct facets. The report obtains this simple graph as the 2-section of the facet-ridge hypergraph: facets incident to a common ridge form a clique. Without the size and positivity conditions, this numerical connectivity claim is not asserted, although the original connectedness conclusion remains. The higher-connectivity statement retains its characteristic-zero restriction and exact polyhedral-hypothesis review needs.
\subsection{Short English statement}
The source asks whether facets of a prime tropicalization are linked through shared ridges. The supplied report attributes an affirmative all-characteristic result to Cartwright--Payne. In characteristic zero it derives a stronger vertex-connectivity and disjoint-facet-chain bound from Maclagan--Yu, subject to lineality and graph-size qualifications.
\subsection{Sources}
\begin{itemize}
\item Local source record: catalogue entry, catalogue entry; canonical \texttt{aim-geometry-notes.json}, \texttt{notes[325]}, item A.6, ``Higher order connectedness of amoebas.'' Title: ``Codimension-one connectivity and robust facet chains.''
\item AIM workshop ``Amoebas and tropical geometry,'' supplied reference \url{https://aimath.org/WWN/amoebas/amoebas.pdf}. Contributors: Manfred Einsiedler, Doug Lind and Rekha Thomas. The verification says AIM HTML resolves OCR damage but gives no separate HTML URL; this draft uses the supplied record without visiting either source.
\item Source background: then-forthcoming work of Einsiedler, Lind and Kapranov. Supplied research attributions: Cartwright--Payne for codimension-one connectedness, Maclagan--Yu for characteristic-zero higher connectivity, and Menger for disjoint graph paths. No full bibliographic entries or exact theorem numbers are supplied.
\item Registry dated 14 August 2026: ``solved,'' ``SOLVED-IN-LITERATURE,'' statement status ``exact''; attempt full solution, rigor ``proved,'' novelty confidence low. The machine-generated, non-peer-reviewed aid requires independent verification before citation.
\end{itemize}
\end{document}
